\bookStart{Fragment of a Lay of Siward}[Brot af Sig·urðar kviða]
\setBookCode{Brot}

\begin{flushright}%
\textbf{Dating} \parencite{Sapp2022}: C10th (0.974)

\textbf{Meter:} \Fornyrdislag%
\end{flushright}

\section{Introduction}

Following the Great Lacuna (\textlink{Lacuna}), the text of \Regius\ picks up in the middle of a heroic poem, namely the present \textbf{Fragment of a Lay of Siward} (\Brot).  This poem (or what survives of it) deals with the murder of Siward and its impact on Byrnhild, his jealous lover and the instigator of his death.

\subsection{Textual preservation and transmission}

\Brot\ is only preserved in \Regius, where it is found on foll.~33r/1--33v/8.  The beginning of the poem has been lost due to the Great Lacuna; see \textlink{Lacuna}.

According to the prose at the end of the poem (P1), \Brot\ originally began with Siward’s death.  Since Siward’s death is implied in st.~4, this indicates that we have not lost too much of the original poem.  On the other hand, it has been speculated that \Brot\ was the original “long” Lay of Siward, contrasting with the “Short Lay of Siward” (\textlink{Sigurdskamma}), in which case we only have a small part of it.  A more likely explanation is perhaps that \textlink{Brot} was preceded by such a poem.

Although part of \Brot\ is lost, the poem clearly served as a source for \VolsungaSaga\ 32, which in fact cites a variant form of st.~4.

\subsection{Meter and stanza}

\Brot\ is composed in \Fornyrdislag\ and consists of 74 complete long-lines distributed across 20 stanzas.  Owing to the loss of the preceding folio, l.~1/1 is defective and consists of only an isolated off-verse.

The typical length of the stanza is 4 lines (17/20: sts.~2--5, 7--16, 18--20) and the poem shows a strong tendency towards the helming-structure.  Other represented lengths are 3 (st.~6) and 2 lines (st.~17).  Stanza 1 originally consisted of at least 2 lines.

\subsection{Summary}

The poem begins \emph{in medias res}.  Hawn asks Guther why he wishes to see Siward dead (1).  Guther answers that Siward has betrayed his marital oaths (and thus their in-lawhood) (2).  Hawn warns Guther that Byrnhild has deceived him; she is jealous of Guthrun since she cannot have Siward (3).  The brothers give their younger brother, Godthorm, the flesh of wolves and snakes in order to strengthen him for the killing of Siward (4).

Guthrun, Siward’s wife, stands outside to greet her brothers and husband as they return; she does not see Siward and asks where he is (5).  Hawn, her brother, answers that they have cut him down with the sword (6).  Byrnhild tells the brothers to enjoy their wealth; if Siward had lived longer he would have taken it all (7).  One of the brothers, probably Hawn again, answers that it would not have been right for Siward to rule their inheritance, since he would have begotten sons to keep it for himself (8).  Byrnhild laughs madly and again tells them to enjoy their lands (9).  Guthrun chastises Byrnhild and then curses her brothers; she tells them that there will be revenge (10).

The poem returns to Siward’s death south of the Rhine; a raven called out above him, telling Guther and Hawn that Attle would soon kill them, that the broken oaths would destroy them (11).  Later the same night, Guther lies awake thinking about this prophecy (12--13).

A little before dawn, Byrnhild awakes; she declares that she will tell of her own grief (14).  All are silent as they hear her weep over the killing which she herself had laughingly instigated (15).  Byrnhild addresses Guther, her husband, with a dream; her speech makes up the rest of the poem.  She dreamt that everything in their hall died; she had a cold bed as Guther rode, fettered, into a host of enemies (16).  So, says she, may all the line of the Nivlings perish; they are oath-breakers (17).  Byrnhild curses her husband for his lack of care for the blood-oaths which he and Siward swore; Siward always put Guther first, and yet he still murdered him (18).  Byrnhild illustrates how much Siward respected Guther by telling her husband how Siward refused to sleep with her while in disguise as him (19).  In disguise as Guther, Siward placed a sword between himself and Byrnhild in their bed.  The outsides of the sword were forged in iron, but its insides were blended with venom---the treachery which ultimately led to their ruin (20).

\sectionline

\newpage

\section{Text}

\bvg\bva%
\speakernote{[Hǫgni:]}%
\Ballnote{The poem appears to begin as Hawn speaks to Guther, who has asked him to help him slay Siward at the behest of Byrnhild.  As usual in heroic tradition (cf.~\textlink{Atlakvida}[6]--12), it is Hawn who is the voice of reason against his rasher brother Guther.}%
\llap{„}\edtrans{\emph{hvat hęfir \alst{S}ig·vǫrðr}}{[What has Siward]}{\Afootnote{om. \Regius}\Bfootnote{A conjectural reconstruction, which however can be held as reasonably secure, owing to the parallelism in the following stanza.}} \hld\ \alst{s}aka unnit &
es þú \alst{f}rǿknan vill \hld\ \alst{f}jǫrvi nę̇ma?“\eva

\bvb\speakernoteb{[Hawn:]}%
“[What has Siward] done for a crime, \\
that thou wilt deprive the brave of life?”\evb\evg


\bvg\bva%
\Ballnote{Although it is most logical that Guther is the one saying this owing to the answer in st.~3, it would perhaps be more fitting for Byrnhild to be the speaker.  In fact, it is an important point of the narrative that Siward has \emph{not} broken his oaths to Guther.  Cf.~the contrasting dialogue \textlink{Sigurdskamma}[13]--19, where Guther is sympathetic towards Siward and aware that he has not broken his oaths.}%
\speakernote{[Gunn·arr:]}%
\llap{„}Mér hęfir \edtext{\alst{S}ig·\emph{vǫ}rðr}{\Afootnote{metr.~emend.~\parencite[135--141]{Sievers1889}; \emph{Sig·urðr} \Regius}} \hld\ \alst{s}ęlda ęiða &
\alst{ęi}ða sęlda \hld\ \alst{a}lla logna; &
þȧ \alst{v}élti hann mik \hld\ es hann \alst{v}esa skyldi &
\alst{a}llra \alst{ęi}ða \hld\ \alst{ęi}nn full-trúi.“\eva

\bvb\speakernoteb{[Guther:]}%
“To me Siward has given oaths: \\
oaths given---all broken! \\
He deceived me when he should have been \\
the one true keeper of all the oaths.”\evb\evg


\bvg\bva%
\speakernote{[Hǫgni:]}%
\llap{„}Þik hęfir \alst{B}ryn·hildr \hld\ \alst{b}ǫl at gørva &
\alst{h}ęiptar \alst{h}vattan \hld\ \alst{h}arm at vinna; &
fyr·man hǫ̇n \alst{G}uð·ru̇nu \hld\ \edtrans{\alst{g}óðra ráða}{her good counsel}{\Bfootnote{Namely her marriage with Siward.}} &
\edtrans{ęn \alst{s}íðan þér \hld\ \alst{s}i̇́n at njóta}{and thereto that thou enjoyest her}{\Bfootnote{I.e., that Guther (rather than Siward) is the one wedded to Byrnhild.}}.“\eva

\bvb\speakernoteb{[Hawn:]}%
“Byrnhild has inflicted a bale upon thee, \\
incited thee to hatred that thou mightst accomplish evil. \\
She begrudges Guthrun her good counsel, \\
and thereto that thou enjoyest her.”\evb\evg


\bvg\bva%
\Ballnote{In order to prepare their younger brother Godthorm for the treacherous killing of Siward, his brothers feed him the meat of wolves, snakes, and dogs. --- A related st.~is cited in \VolsungaMS\ (\VolsungaSaga\ 32); it is sufficiently divergent from the present st.~to be edited on its own as \textlink{Lacuna}[6] above.}%
Sumir \alst{u}lf sviðu, \hld\ sumir \alst{o}rm sniðu, &
sumir \alst{G}oð·þormi \hld\ af \alst{g}ęra dęildu, &
áðr þęir \alst{m}ę́tti \hld\ \alst{m}ęins of lystir &
ȧ \alst{h}orskum \alst{h}al \hld\ \alst{h}ęndr of lęggja.\eva

\bvb Some roasted a wolf; some cut up a snake; \\
some shared hound-flesh with Godthorm, \\
before they might, desiring evil, \\
lay their hands on the wise hero \ken*{= Siward}.\evb\evg


\bvg\bva%
\Ballnote{The killing is implied, but not described explicitly.  In any case, it takes place outside and without Guthrun’s awareness.  This contrasts with \textlink{Sigurdskamma}, \textlink{Gudrunarhvot}, and \textlink{Hamdismal} where Siward is killed in bed next to the unfortunate Guthrun; see note to P1 below.}%
Úti stóð \alst{G}uð·ru̇n \hld\ \alst{G}júka dóttir &
ok hǫ̇n þat \alst{o}rða \hld\ \alst{a}lls fyrst of kvað: &
\llap{„}Hvar ’s nú \edtext{\alst{S}ig·\emph{vǫ}rðr}{\Afootnote{metr.~emend.~\parencite[135--141]{Sievers1889}; \emph{Sig·urðr} \Regius}} \hld\ \alst{s}ęggja dróttinn &
es \alst{f}rę̇ndr mi̇́nir \hld\ \alst{f}yrri ríða?“\eva

\bvb Outside stood Guthrun, Yivick’s daughter, \\
and she this word first of all did say: \\
“Where is now Siward, the lord of men, \\
when my kinsmen ride ahead?”\evb\evg


\bvg\bva%
\edtrans{\alst{Ęi}nn því Hǫgni \hld\ \alst{a}nd·svǫr vęitti}{Alone then Hawn this answer granted}{\Bfootnote{Shared with \textlink{Sigurdskamma}[17]/1, 45/1, which however have \emph{ęinu} (neut. dat. sg.) for \emph{ęinn} (m.~nom. sg.)  The difference is probably a mere spelling error, but it is unclear which form of the line is original.}}: &
\llap{„}\alst{S}undr hǫfum \edtext{\alst{S}ig·\emph{vǫ}rð}{\Afootnote{metr.~emend.~\parencite[135--141]{Sievers1889}; \emph{Sig·urð} \Regius}} \hld\ \alst{s}verði hǫgg\emph{v}inn; &
\alst{g}napir ę́ \alst{g}rár jór \hld\ yfir \alst{g}ram dauðum.“\eva

\bvb Alone, then, Hawn this answer granted: \\
“We have cut Siward asunder by the sword; \\
the grey steed ever neighs over the dead prince.”\evb\evg


\bvg\bva%
Þȧ kvað þat \alst{B}ryn·hildr \hld\ \alst{B}uðla dóttir: &
\llap{„}\alst{V}el skuluð njóta \hld\ \alst{v}ápna ok níu landa; &
\alst{ęi}nn \edtext{myndi}{\Afootnote{\emph{myndiv} \Regius}} \edtext{\alst{S}ig·\emph{vǫ}rðr}{\Afootnote{metr.~emend.~\parencite[135--141]{Sievers1889}; \emph{Sig·urðr} \Regius}} \hld\ \alst{ǫ}llu ráða &
ef hann \alst{l}ęngr \alst{l}ítlu \hld\ \alst{l}ifi hęldi.“\eva

\bvb Then this quoth Byrnhild, Budle’s daughter: \\
“Well shall ye enjoy the weapons and nine lands! \\
Alone would Siward have ruled them all \\
if a little longer he had held his life.”\evb\evg


\bvg\bva\speakernote{[Hǫgni:]}%
\Ballnote{The speaker is not indicated, but since he spoke in st.~6, it is most likely Hawn who responds.}%
\llap{„}Vę́ri⸗a þat \alst{s}ǿmt \hld\ at hann \alst{s}vá réði \lednopb &
\alst{G}júka arfi \hld\ ok \alst{g}ota męngi, &
es hann \edtrans{\alst{f}imm sonu}{five sons}{\Bfootnote{Siward never has five living sons in the heroic tradition (rather, he usually only has one, Syemund).  Instead, Guther is pointing out that Siward, \emph{if} he had lived longer, would have begotten such sons with Guther, and these would eventually have come into conflict with their uncles, the Yivickings.}} \hld\ at \alst{f}olk-rǿði \lednopb &
\alst{g}unnar fu̇sa \hld\ \alst{g}etna hafði.“\eva

\bvb\speakernoteb{[Hawn:]}%
“It would not have been right that he should so have ruled \\
Yivick’s inheritance and the menfolk of the Gots, \\
when to rule the folk he would have begotten \\
five sons eager for war.”\evb\evg


\bvg\bva%
\edtext{Hló þȧ \alst{B}ryn·hildr \hld\ ---\alst{b}ǿr allr dunði--- &
\alst{ęi}nu sinni \hld\ af \alst{ǫ}llum hug:}{\lemma{Hló \dots\ hug ‘Then \dots\ heart’}\Bfootnote{Found near-identically in \textlink{Sigurdskamma}[30]/1--2.}} &
\llap{„}\edtext{Vel skuluð njóta \hld\ landa ok þegna}{\Bfootnote{Alliteration is missing in this line and no obvious emendation can be found; it clearly parallels 7/2 above.}} &
es þér \alst{f}rǿknan gram \hld\ \alst{f}alla létuð.“\eva

\bvb Then Byrnhild laughed---the town all resounded--- \\
a single time out of all her heart: \\
“Well shall ye enjoy the lands and thanes, \\
since ye made the brave prince to fall.”\evb\evg


\bvg\bva%
Þȧ kvað þat \alst{G}uð·ru̇n \hld\ \alst{G}júka dóttir: &
\llap{„}\alst{M}jǫk \alst{m}ę́lir þú \hld\ \alst{m}iklar firnar; &
\edtrans{\alst{g}ramir hafi}{let the fiends have}{\Bfootnote{Apparently a common curse; cf.~\textlink{Harbardsljod}[60].}} \alst{G}unn·ar \hld\ \alst{g}ǫtvað \edtext{\alst{S}ig·\emph{va}rðar}{\Afootnote{metr.~emend.~\parencite[135--141]{Sievers1889}; \emph{Sig·urðar} \Regius}}, &
\alst{h}ęipt-gjarns \alst{h}ugar \hld\ \alst{h}ęfnt skal verða.“\eva

\bvb Then this quoth Guthrun, Yivick’s daughter: \\
“Many great evils hast thou spoken. \\
Let the fiends have Guther, Siward’s burier; \\
for his evil heart there shall be vengeance.”\evb\evg


\bvg\bva%
\edtrans{\alst{S}oltinn varð}{had expired}{\Bfootnote{Lit. ‘became, turned dead’; \emph{soltinn} ‘dead; starved, dejected’ is the past participle of \emph{svelta} ‘die; starve’.  Apart from the present st., \emph{soltinn} is used in three other places in Norse poetry: \textlink{Sigurdskamma}[47]/2 (‘dejected’), 50/3 (‘dead’) and \textlink{GudrunTwo}[10] (‘dejected’).
The use of this participle together with the shared subject matter may indicate a close relationship between these three poems, although the varying depictions of Siward’s death (see P1 below) suggest that \textlink{GudrunTwo} and \Brot\ are more closely related to each other than to \textlink{Sigurdskamma}.}} \edtext{\alst{S}ig·\emph{vǫ}rðr}{\Afootnote{metr.~emend.~\parencite[135--141]{Sievers1889}; \emph{Sig·urðr} \Regius}} \hld\ \alst{s}unnan Ri̇́nar, &
\alst{h}rafn at męiði \hld\ \alst{h}ǫ́tt kallaði: &
\llap{„}\alst{Y}kkr mun \alst{A}tli \hld\ \alst{ę}ggjar rjóða, &
munu \alst{v}íg-ská\emph{a} \hld\ of \alst{v}iða ęiðar.“\eva

\bvb Siward had expired to the south of the Rhine; \\
a raven on a branch loudly called out: \\
“On you two will Attle redden his blades; \\
the oaths will destroy the war-harried men!”\evb\evg


\bvg\bva%
\alst{F}ramm vas kvelda \hld\ \alst{f}jǫlð vas drukkit &
þȧ vas hví-\alst{v}etna \hld\ \alst{v}il-mál talit; &
\alst{s}ofnuðu allir \hld\ es i̇́ \alst{s}ę́ing kvǫ́mu, &
\alst{ęi}nn vakði Gunn·arr \hld\ \alst{ǫ}llum lęngr.\eva

\bvb The evening was late, much was drunk; \\
then were all kinds of pleasant speeches told. \\
All fell asleep when they came into bed; \\
alone did Guther wake longer than all.\evb\evg


\bvg\bva%
\edtext{\alst{F}ót nam at hrǿra \hld\ \alst{f}jǫlð nam at spjalla}{\lemma{Fót \dots\ spjalla ‘His foot \dots\ consider.’}\Bfootnote{Probably a formulaic line-template; cf.~\textlink{Thrymskvida}[1]/3.}} &
\alst{h}itt \alst{h}ęr-glǫtuðr \hld\ \alst{h}yggja téði, &
hvat þęir i̇́ \alst{b}ǫðvi \hld\ \alst{b}áðir sǫgðu &
\alst{h}rafn ęy ok ǫrn \hld\ es þęir \alst{h}ęim riðu.\eva

\bvb His foot he took to move; much he took to consider. \\
Of this the troop-destroyer thought, \\
what in the fray those two had said, \\
the raven always and the eagle, when they rode home.\evb\evg


\bvg\bva%
\Ballnote{Byrnhild awakes just before dawn the day after the drinking; she makes a speech to the whole house.  Her speech takes up the rest of the poem.}%
Vaknaði \alst{B}ryn·hildr \hld\ \alst{B}uðla dóttir &
\alst{d}ís skjǫldunga \hld\ fyr \alst{d}ag lítlu: &
\llap{„}\alst{H}vętið mik eða lętið mik \hld\ ---\edtrans{\alst{h}armr es unninn}{the evil is accomplished}{\Bfootnote{A verbal callback referencing st.~3/2.}}--- &
\alst{s}org at \alst{s}ęgja \hld\ eða \alst{s}vá láta.“\eva

\bvb Byrnhild awoke, Budle’s daughter, \\
the dise of shieldings, a little before day: \\
“Urge me on or restrain me---the evil is accomplished--- \\
to tell my sorrow, or so let it be.”\evb\evg


\bvg\bva%
\Ballnote{Sts.~15--20 are transparently paraphrased in \VolsungaSaga\ 31: \emph{Nú þóttist engi kunna at svara, at Bryn·hildr beiddi þess hlę́jandi, er hón harmaði með gráti.  Þá mę́lti hón: „Þat dreymdi mik, Gunn·arr, at ek átta kalda sę́ng, en þú ríðr í hendr ó·vinum þínum, ok ǫll ę́tt yður mun illa fara, er þér eruð eið-rofa, ok mundir þú þat ó·gløggt, er þit blǫnduðuð blóði saman, Sig·urðr ok þú, er þú rétt hann, ok hefir þú honum allt illu launat þat, er hann gerði vel til þín ok lét þik fremstan vera, ok þá reyndi þat, er hann kom til vár, hvé hann helt sína eiða, at hann lagði okkar í milli it snarp-eggjaða sverð, þat er eitri var hert. [...]}
‘Now nobody thought himself able to respond when Byrnhild now grieved by weeping over that which she had laughed about when she asked for it.
Then she spoke: “I dreamt, Guther, that I had a cold bed, but thou didst ride into the hands of thy enemies,
and all your lineage will fare ill for ye are oath-breakers.
And thou didst not kindly remember how ye two blended your blood together, thou and Siward, when thou betrayed him, and thou hast cruelly repaid him fully for how he was good towards thee and let thee be the foremost,
and when he came to us it was proven how well he kept his oaths,
for he laid between us the sharp-edged sword which was hardened with venom. [...]’}%
\alst{Þ}ǫgðu allir \hld\ við \alst{þ}ví orði &
\alst{f}ár kunni þęim \hld\ \alst{f}ljóða lǫ́tum &
es hǫ̇n \alst{g}rátandi \hld\ \alst{g}ørði⸗sk at sęgja &
þat’s \alst{h}lę́jandi \hld\ \alst{h}ǫlða bęiddi:\eva

\bvb All fell silent upon that word \\
---few men know the manners of women--- \\
when she, weeping, made ready to tell of \\
that which, laughing, she had asked from the heroes:\evb\evg


\bvg\bva%
\llap{„}Hugða’k mér, \alst{G}unn·arr, \hld\ \alst{g}rimmt i̇́ svefni, &
\alst{s}valt allt i̇́ \alst{s}al \hld\ ę́tta’k \alst{s}ę́ing kalda, &
ęn þú \alst{g}ramr riðir \hld\ \alst{g}laums and·vani &
\edtrans{\alst{f}jǫtri \alst{f}atlaðr}{chained in fetters}{\Bfootnote{Cf.~\textlink{Atlakvida}[21].}} \hld\ i̇́ \alst{f}jánda lið.\eva

\bvb “I had a cruel thought, O Guther, in my sleep: \\
everything died in the hall, I had a cold bed, \\
but thou, prince, didst ride without cheerful fellows, \\
chained in fetters, into a troop of foes.\evb\evg


\bvg\bva%
Svá mun \alst{ǫ}ll \alst{y}ður \hld\ \alst{ę́}tt Niflunga &
\alst{a}fli gęngin--- \hld\ eruð \alst{ęi}ð-rofa!\eva

\bvb So will all your line of Nivlings \\
part from power---ye are oath-breakers!\evb\evg


\bvg\bva%
Mant⸗at-tu \alst{G}unn·arr \hld\ til \alst{g}ǫrva þat &
es \edtrans{þit \alst{b}lóði i̇́ spor \hld\ \alst{b}áðir ręnnduð}{your blood in your tracks ye both did drive}{\Bfootnote{Referring to a ritual of blood-brotherhood, wherein the brothers-to-be would spill and mix their blood into their footprints on the ground.  Blood-brotherhood is also mentioned in \textlink{Lokasenna}[9], and the footprint-ritual is described in \textcite{Saxo} 1.6.7: \emph{Siquidem icturi foedus veteres vestigia sua mutui sanguinis aspersione perfundere consueverant, amicitiarum pignus alterni cruoris commercio firmaturi} ‘Now our ancestors, when they meant to strike a pact, would sprinkle their combined blood in their footprints and mingle it, so as to strengthen the pledge of their fellowship.’  For discussion see PCRN History and Structures III:845 ff.}}, &
nú hęfir hǫ̇num þat \alst{a}llt \hld\ \alst{i}llu launat &
es hann \alst{f}ręmstan \edtrans{\emph{þ}ik}{thee}{\Afootnote{sens.~emend.; \emph{sik} ‘himself’ \Regius}\Bfootnote{This reading follows the paraphrase of \VolsungaSaga\ and is undoubtedly correct.  Siward would have Guther as the foremost of his own house and never violated his oaths toward him.}} \hld\ \alst{f}inna vildi.\eva

\bvb Thou didst not recall, Guther, clearly enough \\
that your blood in your tracks ye both did drive. \\
Now hast thou repaid cruelly him for all that, \\
when he would have had \emph{thee} as the foremost man.\evb\evg


\bvg\bva%
\Ballnote{Sts.~19--20 conclude the poem; Byrnhild illustrates Siward’s loyalty towards Guther by telling him how he refused to sleep with her while in disguise.}%
Þȧ \alst{r}ęyndi þat \hld\ es \alst{r}iðit hafði &
\alst{m}óðigr ȧ vit \hld\ \alst{m}i̇́n at biðja, &
\alst{h}vé \alst{h}ęr-glǫtuðr \hld\ \alst{h}afði fyrri &
\alst{ęi}ðum haldit \hld\ við inn \alst{u}nga gram.\eva

\bvb Then it was proven when he had ridden, \\
brave, to visit and ask for my hand, \\
how that destroyer of hosts \ken*{\textsc{warrior} = Siward} had upheld \\
his side of the oaths against the young lord \ken*{= Guther}:\evb\evg


\bvg\bva%
\Ballnote{Cf.~the description in \textlink{Sigurdskamma}[68].}%
\alst{B}ęn-vǫnd of lét \hld\ \alst{b}rugðinn \edtext{g\emph{o}lli}{\Afootnote{\emph{gylli} \Regius}} &
\alst{m}arg-dýrr konungr \hld\ ȧ \alst{m}eðal okkar; &
\alst{ę}ldi vǫ́ru \alst{ę}ggjar \hld\ \alst{ú}tan gǫrvar &
ęn \edtrans{\alst{ęi}tr-dropum}{with drops of venom}{\Bfootnote{Likely reflecting on how Siward’s refusal to sleep with Byrnhild, an act of deep personal loyalty to Guther, still did not dissuade the Yivickings from murdering him.  In fact, it caused it, since it was an act of betrayal against the unwitting Byrnhild, who ultimately incited his death.}} \hld\ \alst{i}nnan fáðar.“\eva

\bvb He placed a wound-wand \ken{sword} wrapped in gold, \\
the much-adorned king, between us two; \\
with fire were its edges on the outside made, \\
but with drops of venom on the inside stained.”\evb\evg


\subsection{Frá dauða Sig·urðar (‘From the Death of Siward’)}

\bpg\bpa%
\Ballnote{This prose passage gives us an interesting example of early source criticism in the treatment of Scandinavian legend.  Siward’s death at the Thing is described above in \textlink{Brot}[4]--6 and later in \textlink{GudrunTwo}[4]--11.  On the other hand, \textlink{Sigurdskamma}[21]--30, \textlink{Gudrunarhvot}[4], and \textlink{Hamdismal}[6], all have him dying in bed alongside Guthrun.}%
\edtext{Hér}{\Afootnote{Introduced by an ornamented 3-line initial in \Regius.}} er sagt í þessi kviðu frá dauða Sig·urðar ok víkr hér svá til sem þeir drę́pi hann úti.  En sumir segja svá at þeir drę́pi hann inni í rekkju sinni sofanda.  En þýðverskir menn segja svá at þeir drę́pi hann úti í skógi, ok svá segir í \edtrans{Guð·rúnar kviðu inni fornu}{the Ancient Lay of Guthrun}{\Bfootnote{Apparently a reference to \textlink{GudrunTwo}[4]--5, although that poem is not given the descriptor \emph{in forna} ‘the Ancient’ in its own title.}} at Sig·urðr ok Gjúka synir hefði til þings riðit þá er hann var drepinn.  En þat segja allir einnig at þeir sviku hann í tryggð ok vógu at hánum liggjanda ok ó·búnum.\epa

\bpb {\huge H}\textsc{ere it has been told} in this lay from the death of Siward, and here it begins at the point when they slew him outside.  But some say that they slew him inside in his bed-chamber, sleeping.  But the Germans say that they slew him outside in the forest, and so it says in the Ancient Lay of Guthrun that Siward and the sons of Yivick had ridden to the Thing when he was slain.  But this they all say in unison, that they betrayed him while he trusted them and struck at him lying down and unarmed.\epb\epg

\sectionline
